\chapter{Introduction}
\label{ch:introduction}

\section{Background}
Smartphone-based fitness tracking applications have become everyday tools for helping individuals monitor aerobic exercise, establish consistent training routines, and quantify energy expenditure. Modern consumer smartphones integrate all necessary telemetry hardware in a single pocketable device: multi-constellation Global Navigation Satellite System (GNSS) receivers, 3-axis MEMS accelerometers and gyroscopes, 3-axis magnetometers, barometric pressure sensors, and operating-system motion co-processors. Because workouts span both open-air environments (parks, urban streets, running tracks) and GPS-denied indoor spaces (gyms, treadmills, covered walkways), a reliable mobile fitness platform must adapt seamlessly between outdoor positioning and indoor inertial dead reckoning without requiring external chest straps or dedicated running watches.

While commercial platforms such as Strava, Nike Run Club, Samsung Health, and Google Fit demonstrate compelling user experiences, their core filtering pipelines, indoor/outdoor transition triggers, step gating heuristics, and calorie scaling factors remain proprietary. Conversely, existing student and open-source mobile fitness projects typically implement isolated fragments---rendering raw polylines directly on a map or polling an unguarded pedometer stream---without addressing urban GNSS drift, operating-system sensor restrictions, stationary false steps, crash-safe local storage, or multi-screen metric synchronization.

The application developed in this thesis is named \textbf{Aetron}. Implemented in Flutter~\cite{flutter2024framework} using a single Dart codebase targeting native Android and iOS devices, Aetron unifies live sensor fusion with offline-first persistence and cloud reconciliation. Location streams are ingested via \texttt{geolocator}~\cite{geolocator2024}, native step counters via \texttt{pedometer}~\cite{pedometer2024} alongside direct triaxial inertial pipelines, reactive state via Riverpod~\cite{riverpod2024state}, local telemetry persistence via Isar~\cite{isar2024database} and SQLite (\texttt{sqflite})~\cite{sqflite2024}, and synchronization plus server-side canonical finalization via Supabase~\cite{supabase2024backend}. Building upon an earlier deterministic pre-thesis baseline~\cite{aetron2026prethesis}, this graduation thesis advances Aetron from heuristic thresholds to peer-reviewed, mathematically grounded algorithms verified against both field trials and open public benchmark datasets.

\section{Problem Statement}
Implementing an accurate, cross-platform fitness tracker exclusively on commodity Android and iOS hardware remains technically demanding because every sensor modality exhibits failure modes across indoor and outdoor settings. GNSS signals degrade severely in urban canyons and fail inside buildings~\cite{merry2019gps,zandbergen2011gps}; hardware pedometers and software peak detectors accumulate false steps when users tap the screen or handle the phone while standing still; operating systems---most notably Apple iOS---strictly sandbox ambient light sensors, cellular radio metrics, and raw satellite carrier-to-noise density; and mobile background execution faces abrupt OS termination.

From a systems and algorithmic perspective, the core problem addressed by this thesis is:
\begin{quote}
\textit{Design, implement, and empirically validate a cross-platform mobile fitness tracking application that accurately captures indoor and outdoor workout sessions, performs robust environment classification across heterogeneous Android and iOS sensor constraints, computes consistent real-time and post-workout telemetry, guarantees crash-resilient offline-first persistence with verifiable cloud reconciliation, and maintains uniform metrics across all application modules.}
\end{quote}

Specifically, the system must resolve five coupled engineering and scientific challenges:
\begin{enumerate}
    \item \textbf{Outdoor GNSS Noise, Multipath Drift, and Estimator Validation:} Smartphone location fixes exhibit stationary wander, multipath jumps, and urban-canyon degradation. A naive fixed-alpha exponential smoother risks freezing real movement at $1$\,Hz, whereas complex adaptive Kalman filters must be empirically verified against independent surveyor-grade ground truth rather than assumed to improve accuracy unconditionally.
    \item \textbf{Cross-Platform Indoor/Outdoor Detection (IOD) under iOS Sandboxing:} Lighting-dependent deep IOD models fail with $0\%$ coverage on iOS because third-party iOS apps cannot query the ambient light sensor or raw GNSS $C/N_0$. Moreover, instantaneous predictions suffer rapid boundary oscillation and activity confounding (misclassifying indoor walking as outdoor locomotion), requiring both light-free multi-sensor modeling and strict safety isolation.
    \item \textbf{Stationary False-Step Rejection and Cadence Stability:} Naive peak-detection and unguarded pedometer streams increment phantom steps during screen interaction or idle sway. Simultaneously, live cadence estimation requires balancing accuracy, second-to-second visual flutter, and stop-decay latency, while static stride constants misestimate varying walking and running speeds.
    \item \textbf{Metric Authority and Physiological Consistency:} Elapsed duration, moving time, pace, distance, cadence, and estimated calories must remain coherent from live recording through post-workout summary, history calendar, analytics charts, detail views, and backend verification---using transparent peer-reviewed lookup standards rather than ad-hoc scaling multipliers.
    \item \textbf{Local-First Durability and Multi-Account Integrity:} Recorded sessions and raw sensor points must survive OS kills, app restarts, connectivity loss, and account switching without losing finalized records or leaking records across signed-in accounts.
\end{enumerate}

\section{Scope and Objectives}
Aetron targets individual users recording outdoor and indoor walking, running, cycling, and gym sessions solely with a smartphone. The application includes multi-modal Flutter/Dart client modules (Home, Activity/Recording, History/Calendar, Analytics, Profile, Goals, and Settings), offline-first Isar/SQLite storage, asynchronous Supabase synchronization, Deno post-processing workers, and a reproducible Python/Dart scientific benchmark suite.

To keep empirical claims honest, the project explicitly excludes: dedicated smartwatch or chest-strap heart-rate integration; clinical indirect calorimetry ($\dot{V}\text{O}_2$ testing)---caloric outputs are explicitly presented as population-average estimates; and production road-network map matching, where an OpenStreetMap/OSRM routing engine is not currently attached (so backend workers transparently report \texttt{routeMatchingAvailable() = false}).

The main objectives of this thesis extend the four pre-thesis goals (\textbf{O1}--\textbf{O4}) and establish formal scientific benchmarking (\textbf{O5}):
\begin{itemize}
    \item \textbf{O1 --- Adaptive Environment Classification and Safe Sensor Routing:} Pair a probabilistic \texttt{EnvironmentClassifier} (sigmoid GNSS likelihoods, Benhamou straightness, causal HMM, and iOS GPS-blackout detection) with cross-platform IOD classifiers (Android 4-channel Model~C-Mobile + Adaptive Transition Engine and iOS-first light-free MagIO + Accelerometer + CoreLocation + Barometer + HMM fusion), enforcing \texttt{isObservationalOnly = true} on experimental models so unverified flips never corrupt workout tracking.
    \item \textbf{O2 --- Research-Grounded Outdoor and Indoor Tracking Pipelines:} Upgrade outdoor filtering to a Joseph-stabilized 4-state kinematic EKF with ZUPT, Phase~1 preprocessing, Phase~2 Mahalanobis NIS gating, T-DBSCAN dwell suppression, and display-only PCHIP outage reconstruction; harden indoor pedometry using causal Oxford WPD and Xu Running detectors with Skog MV-GLRT stance rejection, Brajdic \& Harle \texttt{StepMotionGate}, a 6-second sliding cadence window, and LOSO-validated Renaudin stride estimation.
    \item \textbf{O3 --- Real-Time Session Experience and Caloric Parity:} Maintain live telemetry updates at $\le 1$\,s intervals across a refined Kinetic editorial UI with iOS Live Activities and direct iOS/Android OS notification settings deep-linking, paired with bit-exact 2024 Adult Compendium speed-banded caloric estimation across client and server.
    \item \textbf{O4 --- Crash-Resilient Persistence and Server Reconciliation:} Guarantee local-first Isar persistence with 5-second incremental checkpoints and account-scoped access checks before cloud synchronization, backed by idempotent Supabase PostgreSQL guards and Deno verification workers.
    \item \textbf{O5 --- Empirical Benchmarking, Negative-Result Disclosure, and Automated Verification:} Retain field comparisons against commercial apps alongside rigorous holdout evaluations on open ground-truth datasets (TU~Wien RTK, SnapperGPS~2, Cambridge CIOD, HAR-PMD, STAGA, PedEval, OxWalk, Oxford WPD, Xu Group~3, SLEDataset2, and Fukuchi RBDS), disclose negative results transparently, and maintain green CI (\texttt{0} analyzer issues and full regression test pass).
\end{itemize}

\section{Assumptions and Constraints}
\begin{enumerate}
    \item \textbf{Permissions and Sensor Availability:} Full dual-mode operation assumes the user grants foreground/background location, activity recognition (Android), motion \& fitness (iOS), and optional notification permissions. Notification preferences in settings are synchronized bidirectionally with OS-level authorization state.
    \item \textbf{Platform Sandboxing Constraints:} Apple iOS restricts third-party access to raw ambient illuminance, cellular signal parameters (RSRP/RSRQ), and raw GNSS carrier-phase/$C/N_0$ measurements; therefore, iOS classifiers rely strictly on public \texttt{CMMotionManager}, \texttt{CoreLocation}, and \texttt{CMAltimeter} streams.
    \item \textbf{Biomechanical and Metabolic Boundaries:} Without wearable heart-rate or metabolic-cart equipment, indoor distance from accelerometer/pedometer streams and kilocalorie estimates from 2024 Compendium MET tables represent statistical population approximations rather than individualized clinical measurements.
    \item \textbf{Hardware and Backend Resource Constraints:} Consumer smartphone GNSS accuracy remains constrained by multipath reflections and receiver antennas; backend deployments operate within Supabase database/worker limits without third-party commercial map-matching APIs.
\end{enumerate}

\section{Proposed Solution}
Aetron resolves these challenges through an offline-first, layered architecture combining real-time pure-Dart sensor fusion with transparent backend auditing and rigorous offline verification. At session start, \texttt{WorkoutSessionNotifier} initializes parallel location and motion pipelines alongside the passive \texttt{IodController} observational monitor. Depending on activity modality and real-time classification, verified GNSS updates flow through Phase~1 preprocessing and the 4-state kinematic EKF (or snapshot ZUPT/DBSCAN dwell guards), while indoor or GPS-denied intervals rely on MV-GLRT and \texttt{StepMotionGate}-hardened pedometry with Renaudin adaptive step length.

Every 5 seconds after meaningful progress and at workout completion, states are committed transactionally to Isar before entering the resilient synchronization queue to Supabase and \texttt{deterministic-finalize-worker}. Figure~\ref{fig:proposed_pipeline} illustrates the high-level indoor and outdoor workout data processing pipeline.

\begin{figure}[htbp]
\centering
\noindent\resizebox{0.94\textwidth}{!}{%
\begin{tikzpicture}[
    node distance=0.85cm and 0.95cm,
    box/.style={rectangle, rounded corners=4pt, draw=aetronblue, fill=blue!5, thick, text width=3.0cm, minimum height=1.2cm, align=center, font=\small},
    decision/.style={diamond, aspect=1.7, draw=aetroncyan, fill=cyan!6, thick, text width=2.2cm, align=center, font=\footnotesize, inner sep=1pt},
    obs/.style={rectangle, rounded corners=4pt, dashed, draw=codegray, fill=gray!8, thick, text width=3.85cm, align=center, font=\footnotesize, inner sep=5pt},
    storage/.style={rectangle, rounded corners=4pt, draw=codegreen, fill=green!6, thick, text width=3.0cm, minimum height=1.0cm, align=center, font=\small},
    arrow/.style={-{Stealth[length=2.3mm]}, thick, color=aetronblue}
]
    \node[box] (start) {\textbf{User Starts}\\Workout Session};
    \node[box, right=of start] (acq) {\textbf{Signal Acquisition}\\GNSS, IMU, Mag,\\Baro, Pedometer};
    \node[decision, right=of acq] (env) {\textbf{Tracking}\\Environment?};

    \node[box, above right=0.55cm and 1.55cm of env] (out) {\textbf{Outdoor Pipeline}\\P01 Gate + 4-State EKF\\ZUPT + T-DBSCAN};
    \node[box, below right=0.55cm and 1.55cm of env] (in) {\textbf{Indoor Pipeline}\\MV-GLRT + StepGate\\6s Cadence + Renaudin};

    \node[obs, above=0.65cm of acq] (iod) {\textbf{Observational IOD}\\Android 4-Ch / iOS Light-Free\\(\texttt{isObservationalOnly})};

    \node[box, right=5.65cm of env] (met) {\textbf{Unified Metrics}\\Distance, Moving Time,\\Pace, 2024 Compendium};

    \node[storage, above right=0.55cm and 0.95cm of met] (isar) {\textbf{Local Vault}\\Isar Checkpoints (5s)\\\& SQLite Queue};
    \node[storage, right=0.95cm of met] (ui) {\textbf{Shared Repository}\\Home, History, Stats,\\Calendar, Live Activity};
    \node[storage, below right=0.55cm and 0.95cm of met] (cloud) {\textbf{Supabase Cloud}\\Raw Sync + Deno\\Finalize Workers};

    \draw[arrow] (start) -- (acq);
    \draw[arrow] (acq) -- (env);
    \draw[arrow, dashed, color=codegray] (acq) -- (iod);
    \draw[arrow, dashed, color=codegray] (iod) -| node[pos=0.25, above, font=\scriptsize, color=codegray]{advisory hint} (env.north);

    \draw[arrow] (env.east) |- node[pos=0.76, above=1pt, font=\scriptsize\bfseries]{Outdoor} (out.west);
    \draw[arrow] (env.east) |- node[pos=0.76, below=1pt, font=\scriptsize\bfseries]{Indoor / Gap} (in.west);

    \draw[arrow] (out.east) -| (met.north);
    \draw[arrow] (in.east) -| node[pos=0.25, below, font=\scriptsize]{step fallback} (met.south);

    \draw[arrow] (met) -- (ui);
    \draw[arrow] (met.north east) |- (isar.west);
    \draw[arrow] (isar.east) -- ++(0.35,0) |- (cloud.east);
    \draw[arrow] (met.south east) |- (cloud.west);
\end{tikzpicture}%
}
\caption{End-to-end workout tracking data processing pipeline across outdoor GNSS filtering, indoor inertial pedometry, observational IOD monitoring, local persistence, and cloud reconciliation.}
\label{fig:proposed_pipeline}
\end{figure}

\section{Structure of the Thesis}
This thesis is organized into five chapters:
\begin{itemize}
    \item \textbf{Chapter~1 (Introduction):} Presents the project background, problem statement across indoor and outdoor mobile sensing, scope and objectives (\textbf{O1}--\textbf{O5}), assumptions and constraints, and the proposed end-to-end pipeline.
    \item \textbf{Chapter~2 (Literature Review):} Reviews the mathematical and empirical foundations of outdoor GNSS kinematics, EKF/AEKF and trajectory mining, cross-platform IOD under iOS/Android constraints, MV-GLRT stationary rejection, adaptive stride estimation, 2024 Adult Compendium calorimetry, commercial application baselines, and open research gaps.
    \item \textbf{Chapter~3 (System Design and Implementation):} Details Aetron's Clean Architecture, local Isar/SQLite and Supabase PostgreSQL schemas, probabilistic and multi-sensor IOD engines, 4-state EKF and P01/P02/P03/T-DBSCAN tracking pipelines, MV-GLRT indoor step/cadence system, synchronization and worker reconciliation, cross-screen data consistency, and UX/OS integration.
    \item \textbf{Chapter~4 (Experiment Results and Evaluation):} Reports both live smartphone comparisons (six commercial apps and $800$\,m MAPE) and comprehensive holdout benchmarks on public ground-truth datasets, including disciplined negative-result disclosures, execution profiling, automated regression test counts, and evaluation against objectives \textbf{O1}--\textbf{O5}.
    \item \textbf{Chapter~5 (Conclusion and Future Work):} Summarizes key contributions, critically discloses limitations and threats to validity, and outlines concrete future research steps.
\end{itemize}
